\documentclass[a4paper,11pt]{ctexbook}
\usepackage{finance-style}

\title{\textcolor{themeblue}{\Huge\textbf{金融衍生品与风险管理}}}
\subtitle{\textcolor{gray}{第二卷\quad 共六课}}
\author{}
\date{}

\begin{document}
\maketitle
\thispagestyle{empty}
\tableofcontents
\newpage

% ============================================================
\chapter{远期合约与期货——锁定未来的价格}
% ============================================================

\begin{scenario}
\textbf{面包商和农民如何同时避开价格风险？}

北京的一个面包商担心未来三个月小麦价格上涨——他的生产成本会上升，但面包售价短期内不能频繁调整。同时，河北的一个农民担心三个月后小麦收获时价格会下跌。两人可以签订一份合约——约定三个月后以某个固定价格交易一定数量的小麦。双方都锁定了价格，各自消除了价格风险。这就是远期合约最基本的应用——套期保值（Hedging）。但如果一方违约怎么办？——这就是期货市场存在的理由。
\end{scenario}

\begin{qualitative}
\textbf{远期、期货与其他衍生品的本质}

衍生品（Derivatives）是一类金融合约——其价值"派生"自某种基础资产（商品、股票、债券、利率、汇率等）的价格。它们不创造新的价值，而是\textbf{重新分配风险}。

\textbf{远期（Forwards）vs 期货（Futures）：}
\begin{itemize}
  \item 远期是场外交易（OTC）的一对一合约——条款可定制，但存在对手方违约风险
  \item 期货是交易所交易的标准合约——通过保证金制度和每日盯市（Mark-to-Market）消除违约风险
  \item 远期在到期日结算一次；期货每日结算（逐日盯市）并追加保证金
\end{itemize}

\textbf{关键概念：}套保者（Hedger——锁定价格、消除风险）vs 投机者（Speculator——承担风险并期望从价格变动中获利）vs 套利者（Arbitrageur——利用价格差异获得无风险收益）。套利者的存在确保远期/期货价格与现货价格之间不会出现系统性偏离——"无套利定价"是衍生品定价的基石。
\end{qualitative}

\begin{quantitative}
\textbf{远期价格定价}

无套利条件下，远期价格$F$与现货价格$S_0$的关系：
\[
F = S_0 \times e^{(r + c - q)T}
\]

其中$r$=无风险利率，$c$=存储成本率，$q$=便利收益率。对于金融资产（如股票指数）通常$c=0$、$q$为股息率：
\[
F = S_0 \times e^{(r - q)T}
\]

\textbf{基差（Basis）：}基差 = 现货价格 - 期货价格。基差风险是指基差在套保期间发生不利变动——即使期货价格按预期变动，基差的变化仍可能导致套保不完全。
\[
\text{基差} = S_t - F_t
\]

\textbf{最优对冲比率}
\[
h^* = \rho \times \frac{\sigma_S}{\sigma_F}
\]
其中$\rho$=现货与期货收益率的相关系数，$\sigma_S$和$\sigma_F$分别为现货和期货的标准差。

\textbf{对冲效率：}$h^{*2}$——被对冲掉的风险比例。对能源和农产品，通常在85-95\%之间。
\end{quantitative}

\begin{lessonsummary}
\begin{enumerate}
  \item 远期/期货通过锁定未来价格实现风险转移——套保者消除风险、投机者承担风险
  \item 期货的逐日盯市和保证金制度消除了信用风险，但引入了流动性风险
  \item 远期价格由无套利原则决定——现货价格、持有成本和便利收益
  \item 基差风险是套保不完全的根本原因——最优对冲比率通过历史相关性估计
\end{enumerate}
\end{lessonsummary}

\begin{discuss}
\begin{enumerate}
  \item 中航油事件（2004）：中国航油（新加坡）在石油期货上投机巨亏5.5亿美元。公司CEO陈久霖越权交易（超出董事会批准的套保额度），卖出大量看涨期权收取期权费——本质上是在赌油价不会上涨。当油价从40美元涨到55美元时，保证金追缴导致破产。这个案例的核心教训是什么？（提示：内控缺失 vs 套保与投机的模糊界限）
  \item 2023年，中国推出"中证1000股指期货"和"科创50ETF期权"。为什么一个成熟的股票市场需要衍生品市场？如果没有衍生品，市场会缺少哪些功能？
  \item "套保者"和"投机者"的区分在实际操作中非常模糊。一个航空公司买原油期货——如果它实际消耗原油，是套保；如果它买入的数量远超实际消耗，就是投机。但如何定义"实际消耗"？——这种模糊性为内控失效创造了空间。
\end{enumerate}
\end{discuss}

\begin{exercises}
\begin{enumerate}
  \item 当前黄金现货价格600元/克，无风险年利率3\%，存储成本忽略，一年期黄金远期合约的无套利价格是多少？如果远期报价620元/克，存在什么套利机会？
  \item 一家中国航空公司预计半年后需要购买10万吨航空燃油。燃油现货价格波动率25\%，燃油期货波动率28\%，现货与期货的相关系数0.85。计算最优对冲比率和需要购买的期货合约数量（每张合约100吨）。
  \item 铜现货价格每吨70000元，一年期期货价格72000元，无风险利率3\%，年存储成本1\%。隐含的便利收益率（Convenience Yield）是多少？
\end{enumerate}
\end{exercises}

\xuehaiinline{
\textbf{基差风险的经典案例——美国航空（2008）：}美国航空在2008年前买入大量航空燃油期货来对冲油价上涨风险。但2008年下半年油价从147美元暴跌至33美元，航空公司发现锁定的高价燃油期货产生了巨额账面亏损——虽然他们"实际用油"的成本也下降了，但由于没有采用"期现结合"的会计处理，期货端的亏损全部计入了损益表，导致公司一度濒临破产。这个案例说明：套保的会计处理和会计准则解读可能比套保策略本身更重要。\\
\textbf{中国商品期货市场的全球地位：}中国期货市场在2023年按成交量计为全球最大商品期货市场。上海期货交易所（SHFE）的铜和螺纹钢期货的成交量超过LME；大连商品交易所（DCE）的豆粕和铁矿石期货成交量全球第一；郑州商品交易所（ZCE）的PTA和白糖期货也是全球最大。更重要的是——中国在铁矿石、原油、大豆等关键大宗商品上正在争夺"定价权"——即从作为价格接受者转变为价格制定者。
}

\newpage
% ============================================================
\chapter{互换（Swaps）——利率互换与信用违约互换}
% ============================================================

\begin{scenario}
\textbf{AIG的CDS赌局——1.5万亿美元的"保险"如何击垮了一个保险帝国}

美国国际集团（AIG）在2008年金融危机前卖出了超过5000亿美元名义本金的信用违约互换（CDS）——承诺如果抵押贷款债券违约，AIG来赔。AIG收取了数十亿美元的保费，但几乎为零的保证金——因为评级机构给了AIGAAA评级，交易对手认为AIG"肯定不会违约"。当次贷违约率飙升时，AIG的保证金追缴像雪崩一样涌来——AIG无法支付，美国政府不得不出资1800亿美元救助。问题不在于CDS这个工具本身，而在于AIG对CDS的集中风险敞口（所有鸡蛋在一个篮子里）和完全缺乏抵押品管理。
\end{scenario}

\begin{qualitative}
\textbf{互换的本质}

互换（Swap）是双方约定在未来交换一系列现金流的合约。

\textbf{利率互换（IRS）：}最常见的互换形式——一方支付固定利率、收取浮动利率（反之亦然）。企业用IRS来管理利率风险。例如，一家企业借了浮动利率贷款但预期利率上升——它可以通过IRS把浮动转换为固定，锁定融资成本。

\textbf{信用违约互换（CDS）：}实质上是针对债务违约的"保险"。买方定期支付保费（类似于保险费），卖方在参考实体发生违约时赔付（类似于保险理赔）。
\end{qualitative}

\begin{quantitative}
\textbf{IRS定价}

IRS的固定利率应使互换在签约时价值为零——浮动端与固定端的现值相等：
\[
\text{固定端现值} = C \times \sum_{t=1}^{n} DF_t
\]
\[
\text{浮动端现值} = \sum_{t=1}^{n} F_{t-1,t} \times DF_t
\]
其中$F_{t-1,t}$是远期利率，$DF_t$是折现因子。

固定利率$C$由下式确定：
\[
C = \frac{\sum_{t=1}^{n} F_{t-1,t} \times DF_t}{\sum_{t=1}^{n} DF_t}
\]

\textbf{比较优势理论：}AAA级企业和BBB级企业分别在固定利率市场和浮动利率市场有相对比较优势。通过IRS，双方都可以获得比直接借款更低的融资成本——互换创造了"合成"的融资结构。

\textbf{CDS定价：}保费 ≈ 违约概率 × (1 - 回收率)。违约概率可通过信用利差反推——信用利差 = CDS保费 ≈ 债券收益率 - 无风险利率。
\end{quantitative}

\begin{lessonsummary}
\begin{enumerate}
  \item 利率互换是最基本的利率风险管理工具——将浮动利率暴露转换为固定利率（或反之）
  \item CDS是一种针对信用事件的"保险"——但也常被用作投机工具（无基础资产裸CDS）
  \item 互换市场的核心风险——没有中央清算（场外OTC）、对手方信用风险、集中度风险
  \item AIG的崩溃是互换市场风险失控的典型——集中度+杠杆+监管真空
\end{enumerate}
\end{lessonsummary}

\begin{discuss}
\begin{enumerate}
  \item 2008年后，各国监管机构强制要求标准化的互换通过中央对手方（CCP）清算。这解决了什么问题？又创造了什么新风险？（提示：CCP本身就变成了"大而不能倒"的超级集中点——如果CCP破产呢？）
  \item "裸CDS"（不持有基础资产买入CDS）应该被禁止吗？支持的理由——它是"赌博"而非"保险"；反对的理由——做空者通过CDS发现了次贷的泡沫，如果没有CDS，泡沫可能更大、持续时间更长。
  \item 中国利率互换市场（2019年推出的LPR利率互换）的现状——规模增长缓慢、参与者以银行和券商为主、流动性不足。为什么中国的IRS市场不活跃？是因为利率风险不高，还是因为市场基础设施/参与者专业化程度不足？
\end{enumerate}
\end{discuss}

\begin{exercises}
\begin{enumerate}
  \item 给定以下折现因子和远期利率：6个月DF=0.975, 12个月DF=0.950, 18个月DF=0.925。6个月远期利率F=2.56\%, 12个月远期F=2.63\%。计算一个1.5年期IRS的固定利率（每半年付息一次）。
  \item AAA企业借固定年利率4\%，浮动利率为LIBOR+0.2\%；BBB企业借固定年利率6\%，浮动利率为LIBOR+0.8\%。设计一个IRS，使AAA企业和BBB企业的总融资成本降低，并让双方各获得0.2\%的节省。
  \item 一只5年期公司债的YTM为5.5\%，同期限国债YTM为2.5\%。信用利差300bp。假设回收率40\%，估算隐含的5年累计违约概率。（提示：$-\ln(1 - PD) \times (1 - 回收率) \approx 信用利差$）
\end{enumerate}
\end{exercises}

\xuehaiinline{
\textbf{场外互换市场的"隐性"风险——Wolf of Wall Street的斯坦福版本：}2000-2007年，CDS市场名义本金从1万亿美元扩张到62万亿美元——超过了全球GDP。没有任何中央清算平台、没有任何标准化保证金要求、没有任何集中交易报告。一位交易员戏称这个市场是"相互指向对方的枪支——每一方都声称自己有偿还能力，但如果所有人同时倒下，没有人能站起来"。\\
\textbf{中央对手方的"集中风险"问题：}2008年后，场外衍生品被强制要求通过CCP清算。CCP本身成为一个新的超级系统性风险源——如果CCP出现操作失误（如2020年英国清算所LCH的SwapClear的一次险情）、或保证金设计的"顺周期性"导致危机时的保证金追缴加剧了流动性紧张——CCP的倒下意味着整个互换市场的停摆。CCP的Capital Waterfall（违约损失分摊机制）有清晰的五层结构——但最后一层（CCP自身的剩余资本和来自其他清算会员的评估权）在极端情况下可能不足以覆盖损失。
}

\newpage
% ============================================================
\chapter{期权基础——权利与义务的分离}
% ============================================================

\begin{scenario}
\textbf{Archegos——总收益互换如何摧毁100亿美元}

2021年，家族理财室Archegos Capital Management通过"总收益互换（Total Return Swap）"——一种OTC衍生品——获得了数倍杠杆的中概股和媒体股敞口。Archegos不需要直接持有股份，而是与银行签订TRS合约：银行买入股票，Archegos承担股价变动的全部收益或损失，并支付融资成本。当ViacomCBS等股票暴跌时，Archegos无法追加保证金——多家银行（Credit Suisse, Nomura, Morgan Stanley）被迫平仓，总损失超过100亿美元。这个案例展示了期权的核心特征——杠杆。
\end{scenario}

\begin{qualitative}
\textbf{期权——非线性回报的奇迹}

期权赋予买方"权利"而非"义务"去买卖某种基础资产。这个"权利vs义务"的分离是期权与众不同的核心特征——它创造了不对称的回报结构。

\textbf{看涨期权（Call）：}买方有权利（不是义务）在未来某时间以约定价格买入资产。

\textbf{看跌期权（Put）：}买方有权利（不是义务）在未来某时间以约定价格卖出资产。

\textbf{欧式vs美式：}欧式只能在到期日行权；美式可以在到期日前任何时间行权。

\textbf{内在价值与时间价值：}期权价格 = 内在价值 + 时间价值。时间价值在到期日归零——这就是"Theta"（时间损耗）的来源。
\end{qualitative}

\begin{quantitative}
\textbf{期权到期盈亏}

看涨期权多头（买入Call）：$\max(S_T - K, 0) - C$
看跌期权多头（买入Put）：$\max(K - S_T, 0) - P$

\textbf{看涨-看跌平价关系（Put-Call Parity）——欧式期权的定价锚：}
\[
C + PV(K) = P + S_0
\]

如果这个关系不成立，存在无风险套利机会。这个公式也是BS公式推导的基础。

\textbf{常见组合策略的盈亏图：}
\begin{itemize}
  \item 跨式组合（Straddle）：同时买入Call和Put（同一执行价）——赌大幅波动
  \item 勒式组合（Strangle）：买入价外Call和价外Put——赌大幅波动但成本更低
  \item 蝶式价差（Butterfly）：买入低价Call+卖出两倍中价Call+买入高价Call——赌波动率下降
\end{itemize}
\end{quantitative}

\begin{lessonsummary}
\begin{enumerate}
  \item 期权的核心特征是"权利vs义务"的分离——买方有权利无义务，卖方无权利有义务
  \item 期权的盈亏是非线性的——最大损失有限（保费）、潜在收益无限（Call）/有限（Put）
  \item Put-Call Parity是所有期权定价模型正确性的"试金石"
  \item 期权组合可以构造出几乎任何想得到的收益形态——但没有免费午餐，组合的成本就是波动率
\end{enumerate}
\end{lessonsummary}

\begin{discuss}
\begin{enumerate}
  \item 90\%的期权在到期时一文不值。如果这是真的，为什么还有人买期权？——因为期权提供了不对称的回报（有限损失、潜在暴赚），即便期望值为负，但"彩票效应"（小概率暴赚）吸引了大量投资者。
  \item Robinhood——零佣金期权交易平台——在2021年GameStop事件后受到严厉批评，因为它允许完全没有经验的小额投资者交易高风险期权策略（如裸期权）。期权市场的散户化是一个进步（金融服务民主化）还是一个危机（金融风险散户化）？
  \item "卖出期权"（收取保费、承担潜在大损失）是一种"负期望值的生意"吗？——从统计数据看，长期卖出价外期权的策略确实获得了正的夏普比率，但偶发的极端损失（如"波动率末日"事件）会使多年的积累一次性归零。
\end{enumerate}
\end{discuss}

\begin{exercises}
\begin{enumerate}
  \item 当前股价50元。执行价45元的看涨期权溢价8元，执行价55元的看跌期权溢价3元。画出以下组合的盈亏图：
  \begin{enumerate}
    \item 买入1张Call（行权价45）
    \item 买入1张Put（行权价55）
    \item 同时买入上述Call和Put（跨式组合）
  \end{enumerate}
  \item 股价80元，行权价75元的Call报价8元，行权价85元的Call报价3元，对应的Put价格分别是多少？用Put-Call Parity验证。
  \item 设计一个"零成本"策略——买入一张Call同时卖出一张Put（同一执行价、同一到期日），使得净保费支出为0。这个策略的盈亏结构和直接买入股票相比如何？
\end{enumerate}
\end{exercises}

\xuehaiinline{
\textbf{期权市场的"希腊字母"（Greeks）：}
\begin{itemize}
  \item $\Delta$（Delta）——股价每变动1元，期权价格变动多少。Call的$\Delta$在0到1之间，Put的$\Delta$在-1到0之间。深度实值时$\Delta \to 1$（Call）或$\to -1$（Put）。
  \item $\Gamma$（Gamma）——$\Delta$的变化率。平值期权的Gamma最大——这意味着当股价靠近执行价时，期权的风险暴露变化最快。
  \item $\Theta$（Theta）——时间每流逝一天，期权价格损失多少。期权的买方是Theta的"支付者"（每天亏时间价值），卖方是Theta的"收取者"（每天赚时间价值）。
  \item $V$（Vega）——隐含波动率每变化1\%，期权价格变化多少。Vega在平值、长期期权中最大。
  \item $\rho$（Rho）——利率每变化1\%，期权价格变化多少。Call的Rho为正（利率越高、未来的行权价现值越低、Call越值钱），Put的Rho为负。
\end{itemize}
\textbf{一个经典的"无风险策略"——Delta Hedge：}通过动态调整期权和现货的持仓比例，理论上可以构造一个对短期股价变动免疫的组合。但Gamma的存在使Delta需要不断重新平衡——交易成本和高频调整使得完全对冲在实际中无法实现。LTCM失败的原因之一就是他们认为Delta对冲可以消除价格风险——但他们低估了流动性枯竭时无法以合理价格调整Delta头寸的风险。
}

\newpage
% ============================================================
\chapter{期权定价——Black-Scholes模型}
% ============================================================

\begin{scenario}
\textbf{诺贝尔奖的诅咒——LTCM的诞生与陨落}

Myron Scholes和Fischer Black在1973年发表了期权定价公式——它如此优雅，以至于Scholes和Merton（Robert Merton在同一时期独立推导出了类似结果）在1997年获得了诺贝尔经济学奖（Black在1995年去世，诺奖不追授）。但讽刺的是——获奖后仅一年（1998年），Scholes和Merton参与的对冲基金LTCM就破产了。不是BS公式错了——而是LTCM在50倍杠杆下，用BS公式给非流动性资产定价，当市场波动率（公式中的关键输入）在俄罗斯违约后的几天内飙升了500\%时，LTCM的所有风险模型都失效了。
\end{scenario}

\begin{qualitative}
\textbf{Black-Scholes的核心思想}

BS公式的突破在于——它证明了在连续交易的假设下，可以通过动态对冲完全消除期权的风险，从而推导出期权价格。公式中不包含股票的预期收益——只包含无风险利率。这被称为\textbf{风险中性定价}：在对冲掉所有风险之后，期权的"合适"价格只取决于当前股价、执行价、时间、波动率和无风险利率，和投资者对未来股价走势的看法无关。

\textbf{BS公式的假设：}
\begin{enumerate}
  \item 股价服从几何布朗运动（对数正态分布）
  \item 无交易成本、无税收
  \item 在期权期限内无风险利率恒定
  \item 市场无摩擦、可连续交易
  \item 允许卖空
\end{enumerate}
\end{qualitative}

\begin{quantitative}
\textbf{Black-Scholes定价公式}

欧式看涨期权：
\[
C = S_0 N(d_1) - K e^{-rT} N(d_2)
\]

欧式看跌期权（由Put-Call Parity推导）：
\[
P = K e^{-rT} N(-d_2) - S_0 N(-d_1)
\]

其中：
\[
d_1 = \frac{\ln(S_0/K) + (r + \sigma^2/2)T}{\sigma\sqrt{T}}
\]
\[
d_2 = d_1 - \sigma\sqrt{T}
\]

$N(\cdot)$是标准正态分布的累积分布函数。

\textbf{隐含波动率（Implied Volatility, IV）：}将市场上的期权价格代入BS公式反推出的$\sigma$——它是市场对"未来波动"的一致预期。不同执行价、不同期限的期权通常有不同的IV——这就是"波动率微笑"（Volatility Smile）和"波动率曲面"（Vol Surface）。

\textbf{Vega与Delta对冲：}在BS假设下，Delta对冲可以消除一阶价格风险。但Gamma（$\partial \Delta / \partial S$）和Vega（$\partial C / \partial \sigma$）的风险无法被完全消除——它们是对冲者面临的主要剩余风险。
\end{quantitative}

\begin{lessonsummary}
\begin{enumerate}
  \item BS公式是金融学最优雅的数学成果之一——它将期权定价简化到只需要五个输入参数
  \item BS公式不包含股票的预期收益——这是风险中性定价的奇妙结论
  \item 隐含波动率是市场对"不确定性的定价"——不同执行价的IV差异（波动率微笑）反映了市场对BS假设偏离的补偿
  \item BS公式的实践缺陷——它的假设在现实中几乎全部不成立，但"不断修正的BS"仍然是期权市场的定价基准
\end{enumerate}
\end{lessonsummary}

\begin{discuss}
\begin{enumerate}
  \item 1987年股灾前，美股期权的IV微笑是平的（不同执行价的IV几乎相同）。1987年之后，微笑出现了——SPX价外Put的IV显著高于价外Call的IV。为什么？——因为市场学会了"尾部风险"（极端暴跌风险）在历史上比BS假设更频繁地发生。
  \item LTCM使用的模型中包含了"相关性"的假设——他们假设不同国家债券之间的相关性会保持在历史范围内。当俄罗斯违约后，所有新兴市场债券的相关性同时趋近于1——模型中的"分散化"假设瞬间崩溃。为什么"相关性"这个参数是所有金融模型中最危险的？
  \item 中国50ETF期权的隐含波动率曲面的形状和美国SPX期权有什么不同？——中国期权的IV通常更高（散户为主的市场波动更大）且期限结构更陡峭（短期不确定性更大）。
\end{enumerate}
\end{discuss}

\begin{exercises}
\begin{enumerate}
  \item 给定：$S_0=50$, $K=52$, $T=0.5$年, $r=3\%$, $\sigma=25\%$。计算Call和Put的BS价格（使用标准正态分布表或Python的scipy.stats.norm.cdf）。
  \item 在上题中，如果Call的市场价格是3元，计算隐含波动率（用试错法或数值方法求解）。
  \item 计算Call的$\Delta$、$\Gamma$、$\Theta$、$Vega$、$\rho$：
  \begin{enumerate}
    \item 当$S_0$接近$K$时，哪个Greeks最大？
    \item 当$\sigma$增加时，Call和Put的价格分别怎么变？
    \item 为什么期权的买方被称为"Theta的支付者"？
  \end{enumerate}
\end{enumerate}
\end{exercises}

\xuehaiinline{
\textbf{LTCM的崩溃——来自两位诺贝尔奖得主和华尔街最聪明头脑的活教材：}\\
\textbf{基本面：}LTCM（1994-1998）由John Meriwether（前所罗门兄弟债券交易主管）创立，合伙人包括Myron Scholes（BS公式作者）和Robert Merton（期权定价理论贡献者）。巅峰期资本金约50亿美元，管理资产超过1000亿美元（杠杆约20倍，通过互换等衍生品可达50-100倍名义敞口）。\\
\textbf{策略：}收敛交易（Convergence Trading）——买入被低估的债券、卖出被高估的债券，赌价差回归历史均值。本质上——做空波动率。\\
\textbf{崩溃：}1998年8月17日，俄罗斯政府违约并贬值卢布。全球投资者恐慌性逃向"高质量"资产（买入美国国债、卖出一切有风险的东西）。被LTCM做空的美德国债利差不仅没有收敛，反而大幅扩大——LTCM的模型假设这种相关系数历史上从未出现过。当月LTCM亏损44\%（约20亿美元）。9月初资本金降至6亿美元——在20倍杠杆下只能承受约0.5\%的市场反向变动。美联储组织了14家银行注资36.5亿美元接管LTCM。\\
\textbf{教训：}（1）模型风险——所有模型在市场极端状态下都会失效；（2）杠杆的放大效应——50倍杠杆意味着反向2\%的变动会清零资本金；（3）流动性危机——即使你的头寸理论上是正确的，市场可能在没有流动性的时候就清算了你；（4）"相关性"是所有金融模型中最危险的假设——因为它在最需要分散化的时候失效。
}

\newpage
% ============================================================
\chapter{风险管理——VaR、压力测试与久期匹配}
% ============================================================

\begin{scenario}
\textbf{宝洁的"有毒"互换——一个1.57亿美元的风险管理课}

1994年，宝洁（P\&G）的财务主管签署了一份复杂的利率互换协议——"5s30s利差互换"——他赌美国国债收益率曲线（长端减短端）保持稳定窄幅。结果美联储在1994年加息250bp，收益率曲线剧烈变化——宝洁的互换仓位在几周内亏了1.57亿美元。不是互换本身的问题——问题是P\&G的财务主管没有理解互换的杠杆特征和"非线性"风险。一个世界500强公司的财务部门都无法正确评估利率衍生品的风险——对大多数人来说，VaR和压力测试不是理论工具，而是生存之道。
\end{scenario}

\begin{qualitative}
\textbf{风险管理的四步框架}

\begin{enumerate}
  \item \textbf{识别：}知道自己面对什么风险——市场风险（价格/利率/汇率/波动率）、信用风险（对手方违约）、流动性风险（无法以合理价格平仓）、操作风险（系统故障/人为错误）、模型风险（模型假设错误）
  \item \textbf{度量：}量化风险的规模——VaR、Expected Shortfall、久期缺口
  \item \textbf{管理：}决定如何应对——对冲（转移风险）、分散化（降低集中度）、止损（控制损失）、保险（转移给第三方）
  \item \textbf{监测：}持续监控风险指标的变化——设置限额、触发条件和应急预案
\end{enumerate}
\end{qualitative}

\begin{quantitative}
\textbf{VaR（Value at Risk）}

VaR回答的问题是："在$1-\alpha$的置信水平下，未来$T$天内最大可能损失是多少？"

\textbf{三种计算方法：}
\begin{itemize}
  \item \textbf{方差-协方差法：}假设收益正态分布。VaR$_{95\%} = \mu - 1.645\sigma$。简单但无法捕捉尾部风险。
  \item \textbf{历史模拟法：}用过去N天的收益率分布来确定分位点。不依赖分布假设，但假设"历史会重演"。
  \item \textbf{蒙特卡洛模拟法：}模拟大量情景下的价格路径。灵活但计算量大、依赖随机过程假设。
\end{itemize}

\textbf{VaR的致命缺陷：}VaR不告诉你"当损失超过VaR时，损失有多大"——它只告诉你一个阈值。两个投资组合可以有相同的VaR$_{95\%}=1000万美元，但一个的超额损失期望（Expected Shortfall, ES）可能是1200万，另一个可能是5000万。

\textbf{ES（Expected Shortfall）——VaR的改进版：}
\[
ES_\alpha = E[L | L > VaR_\alpha]
\]

巴塞尔协议III（2019年修订版）已将VaR替换为ES作为市场风险资本计量的标准方法。

\textbf{银行的久期缺口管理：}
\[
\text{净资产变动} \approx -\text{久期缺口} \times \Delta y \times \text{总资产}
\]
其中久期缺口 = 资产久期 - (负债/资产) × 负债久期。
\end{quantitative}

\begin{lessonsummary}
\begin{enumerate}
  \item 风险管理是识别→度量→管理→监测的循环过程——不是一次性工作
  \item VaR是最广泛使用的风险度量工具，但有根本缺陷——不度量尾部损失的大小
  \item ES（Expected Shortfall）取代VaR成为新的监管标准——它考虑了尾部风险
  \item 久期缺口管理是商业银行资产负债管理的核心——久期错配是SVB破产的根本原因
\end{enumerate}
\end{lessonsummary}

\begin{discuss}
\begin{enumerate}
  \item 2008年金融危机前，华尔街各大银行的VaR模型都显示"99\%安全的"——每天99\%的概率损失不超过某个数。但模型中没有包含流动性风险和相关性突变——当所有资产同时暴跌、相关性趋近于1时，VaR完全失去了意义。这是"模型错误"还是"模型被误用"？
  \item "压力测试"（Stress Testing）的优势在于它的灵活性——你可以设计任何你想测试的情景（市场暴跌30\%、利率上升300bp、GDP收缩5\%）。但它有两个问题：第一，你无法预知未来会发生的"黑天鹅"是什么样子的；第二，即使你设计了一个"前所未有的"压力情景，也没有人知道市场在这个情景下的真实反应。压力测试是"有用的自我安慰"还是"真正有效的风险管理工具"？
  \item 中国商业银行的"拨备覆盖率"和"资本充足率"在2024年分别约180\%和15\%——远高于监管最低要求和国际同行。高资本缓冲是否意味着中国银行体系比西方更"安全"？还是说这些数字本身——基于特定的不良贷款认定标准——可能低估了真实风险？
\end{enumerate}
\end{discuss}

\begin{exercises}
\begin{enumerate}
  \item 给定一个投资组合过去500个交易日的每日收益率序列，手动或使用Python：
  \begin{enumerate}
    \item 在95\%置信水平下计算VaR（历史模拟法）
    \item 计算95\%置信水平下的Expected Shortfall
    \item 如果假设收益服从正态分布，VaR和ES分别是多少？（用历史和正态两种方法对比）
  \end{enumerate}
  \item 一家银行的资产40亿，负债36亿。资产久期4.5年，负债久期1.2年。利率上升100bp时，银行净资产变动多少？这个变动占净资产的百分比是多少？如果利率上升300bp呢？（回想SVB——当久期缺口超过3年、利率上升500bp、大多数存款是无保险的——灾难的配方。）
  \item 设计一个综合压力测试：假设市场暴跌20\%、信用利差扩大200bp、同时流动性下降（买卖价差从0.1\%扩大到1\%）。对一个"60\%股票+40\%债券"的基准组合，压力情景下的总损失是多少？
\end{enumerate}
\end{exercises}

\xuehaiinline{
\textbf{久期缺口的"死亡螺旋"——SVB案例的技术解剖：}
\begin{itemize}
  \item SVB总资产约2118亿美元，其中约1200亿是持有至到期（HTM）的国债和MBS——久期约6.2年
  \item 负债端主要是活期存款（久期接近0）——虽然少数存款有固定期限，但SVB的存款中94\%超过FDIC保险上限（25万美元），这些"热钱"可以在几分钟内撤走
  \item 久期缺口 ≈ 6.2年 - (2118/2118) × 0 ≈ 6.2年。但更精确的分析应考虑负债端大额存款的"有效久期"——因为储户可以随时取款，负债久期实际上比账面久期更短——缺口更大
  \item 美联储加息500bp导致的市值损失 ≈ 6.2 × 5\% × 2118 ≈ 658亿美元——这远超SVB的总股本（约160亿美元）
  \item SVB倒下的钥匙不在于"资产是否值得"（它的资产是高质量的国债）——而在于这些高质量资产的按市值计价损失超过了其资本金，而存款人同时撤资迫使SVB在不利条件下变现HTM资产
\end{itemize}
\textbf{Basel III/IV的流动性覆盖率（LCR）和净稳定资金比率（NSFR）：}如果SVB严格执行这些指标，它会需要持有更多的高质量流动资产（HQLA）和更稳定的长期资金。但SVB在美国中小银行中属于受监管较宽松的第四类机构（资产规模2500亿以下）——这是监管的漏洞。
}

\newpage
% ============================================================
\chapter{结构型产品与资产证券化}
% ============================================================

\begin{scenario}
\textbf{"杀死华尔街的数学公式"——Gaussian Copula如何制造了次贷危机}

2005年，David Li发表了一篇论文——《On Default Correlation: A Copula Function Approach》。他提出了一个优雅的数学框架来估计多个债务人的联合违约概率——Gaussian Copula（正态连接函数）。华尔街迅速将其应用于CDO（担保债务凭证）的定价。这个公式让复杂的CDO分层变成了一件简单的事——只需要输入一个相关性参数$\rho$。但问题在于：Li假设相关性是恒定的——而在危机中所有相关性都趋近于1。2008年次贷崩盘时，这个"杀死华尔街的公式"被指责为金融危机的技术推手。
\end{scenario}

\begin{qualitative}
\textbf{资产证券化的流程}

\begin{enumerate}
  \item \textbf{贷款池：}银行发放了大量房贷/车贷/信用卡贷款
  \item \textbf{SPV（特殊目的载体）：}将这些贷款卖给一个独立的法律实体（SPV）
  \item \textbf{分级（Tranching）：}SPV发行不同优先级的证券——Senior（AAA）、Mezzanine（BBB、A）、Equity（未评级）
  \item \textbf{信用增级：}Senior层级获得Mezzanine和Equity层的"信用保护"——即使贷款违约，也先损失Equity、再Mezzanine、最后Senior
  \item \textbf{销售：}将不同层级的证券卖给不同风险偏好的投资者
\end{enumerate}

\textbf{关键——分层结构让原本风险很高的次级贷款，通过"信用增级"变成了AAA评级的证券。}这个转化的"黑魔法"就是Gaussian Copula。

\textbf{CDO的"二次加工"：}合成CDO（Synthetic CDO）甚至不持有任何基础贷款——它只包含一篮子CDS合约。再将这些CDS的"风险"再次分层和打包。结果——基础资产是什么已经完全不重要了，投资者甚至不知道自己买的是什么。
\end{qualitative}

\begin{quantitative}
\textbf{Gaussian Copula模型}

单个资产在时间$T$前的违约概率$Q_i(T)$：
\[
Q_i(T) = 1 - e^{-\lambda_i T}
\]
其中$\lambda_i$为违约强度。

Gaussian Copula定义联合违约分布：
\[
P(\tau_1 < T, \tau_2 < T, \ldots, \tau_n < T) = \Phi_n(\Phi^{-1}(Q_1(T)), \Phi^{-1}(Q_2(T)), \ldots, \Phi^{-1}(Q_n(T)); \Sigma)
\]

其中$\Phi_n$为$n$元正态分布的CDF。当$n$很大时，模型只需要输入每个资产各自的违约概率和一个共同的相关性参数$\rho$。这就是为什么华尔街如此热爱这个模型——它把极其复杂的问题简化到了可操作的水平。

\textbf{CDO分级的预期损失计算：}
\begin{enumerate}
  \item 模拟资产池中每个贷款的违约时间
  \item 计算各层级吸收的累计损失额
  \item 重复上万次蒙卡模拟，得到各层级的预期损失和评级
\end{enumerate}
\end{quantitative}

\begin{lessonsummary}
\begin{enumerate}
  \item 资产证券化通过结构分层实现了信用风险的"重新分配"——但并没有消灭风险
  \item Gaussian Copula的致命假设——违约相关性恒定——在危机中完全失效
  \item "合成CDO"让风险变得不透明——投资者不知道底层资产是什么
  \item 2008年后结构化金融被妖魔化——但证券化本身不是问题，问题在于劣质基础资产+过度分层+评级腐败+零透明度
\end{enumerate}
\end{lessonsummary}

\begin{discuss}
\begin{enumerate}
  \item "评级机构在CDO评级中扮演了什么角色？"——CDO发行者（投行）雇佣评级机构（标准普尔、穆迪、惠誉）来为CDO分层评级。存在根本性的利益冲突——评级机构从投行收取评级费，有动力给出更高的评级来争取业务。2008年后Dodd-Frank法案规定评级机构必须接受更严格的监管——但利益冲突的结构性根源并未消除。
  \item 中国的信贷资产证券化（CLO）和RMBS市场在2024年存量约4万亿。为什么中国没有出现次贷危机？原因：基础资产质量较好（中国的按揭贷款首付比例不低于20-30\%——是有效的保护垫）、没有"次级贷款"产品、银行持有大量贷款而非出表（保留"皮肤参与"）。
  \item ABS（资产支持证券）和CDO的本质区别——ABS通常是单层结构（风险不经过复杂分层），CDO是多层结构（经过多次分层和再打包）。多层结构的风险在哪里？——每一次分层都增加了模型假设的数量（相关性、回收率、提前还款率……），每一个假设都可能出错。当你依赖5-10个假设时，全部同时成立的概率接近零。
\end{enumerate}
\end{discuss}

\begin{exercises}
\begin{enumerate}
  \item 简化CDO模拟：假设资产池包含100笔贷款，每笔100万元。每笔贷款的违约概率独立为5\%，但违约事件之间的相关性为20\%（Gaussian Copula设定）。用蒙卡模拟（至少10000次）：
  \begin{enumerate}
    \item 资产池总损失的分布——平均损失、95\%分位损失
    \item 如果资产池被分为Senior（0-80\%损失由Equity和Mezzanine吸收）和Equity（0-5\%）+ Mezzanine（5-20\%）+ Senior（20-100\%），各层级的预期损失率是多少？
  \end{enumerate}
  \item 将违约相关性从0.2提高到0.6，重新计算。更高相关性对Senior层级的保护效果有什么影响？
  \item 案例研究：中国的信贷资产证券化（CRMW——信用风险缓释凭证）。CRMW和CDS的关键区别是什么？为什么中国监管层更鼓励CRMW而不是CDS？
\end{enumerate}
\end{exercises}

\xuehaiinline{
\textbf{Gaussian Copula的完整批判：}
\begin{enumerate}
  \item \textbf{相关性的不稳定性：}Li的模型假设相关性是恒定的——但实证研究表明，相关性在市场压力时期系统性上升。Longin \& Solnik（2001）发现了国际股市在熊市中的相关系数显著大于牛市中。这就意味着——正是你最需要分散化（保护CDO的Senior层）的时候，相关性的上升使保护效果最差。
  \item \textbf{"
同一个$\rho$给所有人"：}Gaussian Copula在华尔街被如此广泛采用的原因之一是它的简洁——给定一个$\rho$就可以给整个CDO分层。"你的参数是多少？"——"0.15。"——"好，就是它了。"没有人质疑$\rho=0.15$是否合理地反映了基础资产的真实相关性结构。
  \item \textbf{尾部独立性：}Gaussian Copula的一个关键数学特性——它在尾部（极端情况下）假设渐进独立。两个极端事件同时发生的概率在Gaussian Copula中被严重低估。而真实数据（如金融危机中多个资产同时违约）表现出显著的"尾部相依性"——这正是Copula的致命伤。
\end{enumerate}
\textbf{中国ABS市场的现状与展望：}到2024年，中国ABS市场存量约4万亿人民币，以信贷ABS（企业贷款和个人消费贷）和RMBS（住房按揭贷款支持证券）为主。发展瓶颈：（1）基础资产池的信息披露不充分——投资者无法独立评估贷款质量；（2）发起人自留比例不足（中国规定发起人须自留5\%——理论上提供了"皮肤参与"，但在实践中可能通过其他方式转移风险）；（3）法律上的"真实出售"认定不够清晰——如果原始人破产，SPV是否真正独立于破产财产？中国的资产证券化要继续发展，必须在上述三个方向有实质性改进。
}

\end{document}
